\documentclass{article}
\usepackage[margin=2cm]{geometry}
\usepackage{amsmath}
\usepackage{graphicx}
\usepackage{booktabs}
\usepackage[utf8]{inputenc}
\begin{document}
\subsection*{0.1 Mass deformations}
In the following, we consider the coset with \(n=1\), i.e. a single vector multiplet. The coset representative is expressed in terms of four scalars \(\phi^{i}, i=0,1,2,3\) via
\[
L=\prod_{i=0}^{3} e^{\phi^{i} K^{i}}
\]
where \(K^{i}\) are the non compact generators of \(S O(4,1)\); see for details. Note that \(\phi^{0}\) is an \(S U(2)_{R}\) singlet, while the other three scalars \(\phi^{r}\) form an \(S U(2)_{R}\) triplet. The scalar potential for this specific case can be obtained from and takes the following form
\[
\begin{aligned}
V\left(\sigma, \phi^{i}\right)= & -g^{2} e^{2 \sigma}+\frac{1}{8} m e^{-6 \sigma}\left[-32 g e^{4 \sigma} \cosh \phi^{0} \cosh \phi^{1} \cosh \phi^{2} \cosh \phi^{3}+8 m \cosh ^{2} \phi^{0}\right. \\
& \left.+m \sinh ^{2} \phi^{0}\left(-6+8 \cosh ^{2} \phi^{1} \cosh ^{2} \phi^{2} \cosh \left(2 \phi^{3}\right)+\cosh \left(2\left(\phi^{1}-\phi^{2}\right)\right)\right.\right. \\
& \left.+\cosh \left(2\left(\phi^{1}+\phi^{2}\right)\right)+2 \cosh \left(2 \phi^{1}\right)+2 \cosh \left(2 \phi^{2}\right)\right)\right] \\
& \left.\left.\left(\right.\right.\right)\left.\left.\left.\left.\left.\left.\left.\left.\left.\left.\left.2\left(\phi^{1}+\phi^{2}\right)\right)\right)+2 \cosh \left(2 \phi^{1}\right)+\left.2 \cosh \left(2 \phi^{2}\right)\right)\right]\right.\right.\right.\right.\right.\right.\right.
\end{aligned}
\]
The supersymmetric \(\mathrm{AdS}_{6}\) vacuum is given by setting \(g=3 m\) and setting all scalars to vanish. The masses of the linearized scalar fluctuation around the AdS vacuum determine the dimensions of the dual scalar operators in the SCFT via
\[
m^{2} l^{2}=\Delta(\Delta-5)
\]
where \(l\) is the curvature radius of the \(\mathrm{AdS}_{6}\) vacuum. For the scalars at hand, one finds
\[
m_{\sigma}^{2} l^{2}=-6 \quad m_{\phi 0}^{2} l^{2}=-4 \quad m_{\phi r}^{2} l^{2}=-6, \quad r=1,2,3
\]
Hence the dimensions of the dual operators are
\[
\Delta_{\mathcal{O}_{\sigma}}=3, \quad \Delta_{\mathcal{O}_{\phi 0}}=4, \quad \Delta_{\mathcal{O}_{\phi^{r}}}=3, \quad r=1,2,3
\]
In these CFT operators were expressed in terms of free hypermultiplets (i.e. the singleton sector). The case of \(n=1\) corresponds to having a single free hypermultiplet, consisting of four real scalars \(q_{A}^{t}\) and two symplectic Majorana spinors \(\psi^{l}\). Here \(I=1,2\) is the \(S U(2)_{R}\) R-symmetry index and \(A=1,2\) is the \(S U(2)\) flavor symmetry index. The gauge invariant operators appearing in are related to these fundamental fields as follows,
\[
\mathcal{O}_{\sigma}=\left(q^{*}\right)^{A} I q^{I}{ }_{A}, \quad \mathcal{O}_{\phi 0}=\bar{\psi}_{I} \psi^{I}, \quad \mathcal{O}_{\phi^{r}}=\left(q^{*}\right)^{A} I\left(\sigma^{r}\right)_{A}^{B} q^{I}{ }_{B}^{B}, \quad r=1,2,3
\]
Note that the first two operators correspond to mass terms for the scalars and fermions, respectively, in the hypermultiplet. The third operator is a triplet with respect to the
\end{document}
